\chapter{Conclusão}
\label{cap:conclusao}

Este trabalho abordou o uso de tecnologias conversacionais no apoio ao controle financeiro pessoal, tendo como ponto de partida as dificuldades de manter o registro e o acompanhamento das movimentações no cotidiano. A partir desse problema, foi desenvolvido o GranaAI, que reúne um assistente integrado ao WhatsApp e uma plataforma web para organização e consulta das informações financeiras.

A pesquisa exploratória, composta por 35 respostas, contribuiu para identificar necessidades e orientar a definição dos requisitos. A busca por praticidade no registro e a utilização de meios familiares de comunicação fundamentaram a adoção do WhatsApp como canal de interação. O levantamento também apoiou a escolha de recursos voltados ao acompanhamento de gastos e objetivos financeiros, estabelecendo uma relação entre a investigação inicial e as decisões de desenvolvimento.

Com base nesses requisitos, foram implementados o registro e a consulta de movimentações por mensagens, juntamente com os recursos de transações, categorias, metas, assinaturas e relatórios da plataforma web. A integração entre essas interfaces permitiu vincular as operações realizadas na conversa às informações apresentadas no dashboard. Dessa forma, o desenvolvimento reuniu em uma mesma aplicação o registro conversacional e o gerenciamento dos dados financeiros.

A contribuição técnica do projeto está nessa articulação entre duas formas de acesso às informações: a conversa para registrar e consultar movimentações e a plataforma web para organizá-las e acompanhá-las de forma consolidada. O GranaAI representa, portanto, uma aplicação das interfaces conversacionais ao controle financeiro pessoal, em que o uso de mensagens passa a fazer parte do próprio fluxo de registro e consulta do sistema.
